\section{Week 1 Transcript Seam}
\label{sec:week1-retro}

Week 1 established the first extracted cross-module theorem on the challenge
suffix path:
\[
\take_{32}(\mu_{\mathsf{sign}}^{64}) = \mu_{\mathsf{verify}}^{32}
\Longrightarrow
\mathsf{equal\_challenge\_suffix}.
\]

\paragraph{Primary claim IDs.}
\begin{itemize}
  \item \claimid{MU32-ABSORB}: \status{PROVED}.
  \item \claimid{CHAL-LIST}: \status{PROVED}.
  \item \claimid{CHAL-CALL}: \status{PARTIAL}.
\end{itemize}

\paragraph{Pinned procedures and theories.}
\begin{itemize}
  \item Sign procedure: \procname{__sign\_challenge\_shake256\_absorb\_mu32}.
  \item Verify procedure: \procname{__verify\_shake256\_absorb\_mu32}.
  \item EasyCrypt theorem:
        \breakablesame{ExtractedChallengeAbsorb.sign_verify_mu32_absorb_from_pos64}.
  \item EasyCrypt theorem:
        \breakablesame{TranscriptBytes.challenge_input_eq_components}.
\end{itemize}

\paragraph{Proof strategy.}
The extracted theorem fixed sponge position 64 and required an explicit
\texttt{mu32\_prefix} predicate over generated arrays. A separate witness lemma
showed that this premise is satisfiable, which ruled out one vacuity risk but
did not yet connect the premise to actual KeyGen/Sign/Verify call paths.

\paragraph{Verification evidence.}
Week 1 used fresh \texttt{-no-eco} compilation and regenerated focused
extractions. The operational source of truth remains the Markdown ledger and
the verifier logs under \pathname{haetae-topdown-easycrypt/logs/}.
